\documentclass[11pt]{article}
\usepackage[margin=0.9in]{geometry}
\usepackage{amsmath,amssymb,amsthm,hyperref}
\newtheorem{theorem}{Theorem}
\newtheorem{lemma}{Lemma}
\newcommand{\RR}{\mathbb R}
\title{Coinciding symmetries determine a homogeneous mean}
\author{Principia Math In Progress}
\date{27 September 2026}
\begin{document}
\maketitle
\begin{abstract}
We prove Hypothesis 1 in the final version of Mihokovi\'c's
\emph{Coinciding Mean of the Two Symmetries on the Set of Mean Functions}.
In fact the two symmetry operations, together with only a second-order
diagonal expansion, determine the entire symmetric homogeneous mean.
The result applies to every real parameter, as required by the final
article, and does not assume differentiability away from the diagonal.
\end{abstract}

\section*{Task 1: the current source question}
A strict symmetric mean is a function $M:(0,\infty)^2\to(0,\infty)$
with $M(a,a)=a$, $M(a,b)=M(b,a)$ and
$a<M(a,b)<b$ whenever $a<b$. It is homogeneous when
$M(\lambda a,\lambda b)=\lambda M(a,b)$ for $\lambda>0$.
For another strict symmetric mean $N$, the group symmetry in
\cite[Definition 1 and (4)]{M} is, off the diagonal,
\begin{equation}\label{S}
 S_M(N)(a,b)=
 \frac{a(u-a)(m-b)^2-b(m-a)^2(u-b)}
 {(u-a)(m-b)^2-(m-a)^2(u-b)},\quad
 m=M(a,b),\quad u=N(a,b).
\end{equation}
On the diagonal its value is $a$.
The functional symmetry is characterized by
\begin{equation}\label{F}
 M\bigl(N(a,b),\sigma_M(N)(a,b)\bigr)=M(a,b).
\end{equation}
The question concerns means for which these symmetries coincide for
every $N$. Hypothesis 1 on page 15 of the final article asks whether
the asymptotic coefficients of every such homogeneous mean agree
with those of $L_c$ for some $c\in\RR$.

The MathDB wording and arXiv manuscript restrict $c$ to $[-1,\infty)$.
The final article corrects this to $\RR$: its Corollary 1 and Theorem 3
also use all real $c$. We solve this current formulation. The old
restricted-range statement is already superseded and is not the basis
for new credit here.

\begin{theorem}\label{main}
Suppose $M$ is a strict symmetric homogeneous mean, the identity
\eqref{F} holds with $\sigma_M(N)=S_M(N)$ for every strict symmetric
mean $N$, and there are real $\beta,c$ such that
\begin{equation}\label{exp}
 M(x-1,x+1)=x+\beta+\frac{c}{x}+o(x^{-1})\qquad(x\to\infty).
\end{equation}
Then $\beta=0$ and $M=L_c$ everywhere, where
\begin{equation}\label{Lc}
 L_c(a,b)=
 \begin{cases}
 \displaystyle\frac{(1+2c)(a+b)+
 \sqrt{(a+b)^2+4c(1+c)(b-a)^2}}{4(1+c)},&c\ne-1,\\[6pt]
 \displaystyle\frac{2ab}{a+b},&c=-1.
 \end{cases}
\end{equation}
Conversely, for every real $c$, $L_c$ is a strict symmetric homogeneous
mean and its two symmetries coincide for every $N$.
\end{theorem}

The source's full asymptotic expansion implies \eqref{exp}. Thus the
theorem proves more than coefficient agreement, under weaker expansion
requirements. The possible constant $\beta$ is included to avoid
assuming in advance that symmetry alone excludes an absolute-value
term at the diagonal.

\section*{The local perturbation}
Fix $0<a<b$, let $m=M(a,b)$, and write $d=m-a>0$, $e=b-m>0$.
For $u=m+h$ let $v$ denote the right side of \eqref{S}. Its denominator
is positive for $a<u<b$: it is $(u-a)e^2+d^2(b-u)$.
Direct simplification gives
\begin{equation}\label{reflection}
 v=m-\frac{h}{1+\alpha h},\qquad
 \alpha=\frac{e-d}{de}=\frac{a+b-2m}{(m-a)(b-m)}.
\end{equation}
In particular, as $h\to0$,
\[
 v=m-h+\alpha h^2+O(h^3),\qquad
 s:=\frac{u+v}{2}=m+\frac\alpha2h^2+O(h^3),\qquad
 r:=\frac{|u-v|}{2}=|h|+O(h^2).
\]

These perturbations are realized by admissible means. Choose $R=A$,
the arithmetic mean, if $A(a,b)\ne m$; otherwise choose the harmonic
mean $R=H$, since $H(a,b)<A(a,b)$ off the diagonal. For
$0<\lambda<1$, the mean $N_\lambda=(1-\lambda)M+\lambda R$ is strict,
symmetric and homogeneous, and
$h=\lambda(R(a,b)-m)$ is nonzero and tends to zero from one side.
The symmetry hypothesis therefore gives $M(u,v)=m$ along this limit.
If the tested means are also required to possess full asymptotic
expansions, these convex combinations do whenever $M$ does.

By homogeneity, \eqref{exp} is equivalent to
\[
 M(1-z,1+z)=1+\beta z+c z^2+o(z^2)\qquad(z\downarrow0).
\]
Since $M$ is symmetric, we may apply it with $z=r/s\ge0$ and obtain
\begin{equation}\label{local}
 m=M(u,v)=s+\beta r+c\frac{r^2}{s}+o(r^2/s).
\end{equation}
After subtracting $m$ and dividing by $|h|$, the limit in
\eqref{local} gives $\beta=0$. Dividing next by $h^2$ yields
\[
 \frac\alpha2+\frac{c}{m}=0.
\]
Substituting \eqref{reflection} and multiplying by the nonzero
quantity $m(m-a)(b-m)$ shows that
\begin{equation}\label{quad}
 Q_c(a,b;m):=2(1+c)m^2-(1+2c)(a+b)m+2cab=0.
\end{equation}
All limits were one-sided along admissible means; no differentiation
of $M$ at $(a,b)$ was used.

\section*{The unique interior root and the converse}
For every real $c$,
\[
 Q_c(a,b;a)=a(a-b)<0,\qquad Q_c(a,b;b)=b(b-a)>0.
\]
A polynomial of degree at most two with these endpoint signs has
exactly one root in $(a,b)$. Indeed, existence follows by continuity;
two distinct interior roots would force the endpoint signs to agree,
and a double root does not change sign. For $c=-1$, the polynomial
is linear and the root is $2ab/(a+b)$.
For $c\ne-1$ its discriminant is
\[
 (a+b)^2+4c(1+c)(b-a)^2
 =(1+2c)^2(a-b)^2+4ab>0.
\]
The root where the polynomial crosses from negative to positive is
exactly the plus-square-root expression in \eqref{Lc}, for either
sign of $1+c$. This proves $M=L_c$ and also proves strict internality
of $L_c$. Its symmetry, homogeneity and diagonal identity follow
directly from the formula, or from uniqueness of the interior root.

For completeness we prove symmetry coincidence on this entire real
parameter range. Fix $a<b$, put $m=L_c(a,b)$, $u=N(a,b)$, and let
$v$ be \eqref{S}. The positive odds relation
\[
 \frac{v-a}{b-v}=
 \left(\frac{m-a}{b-m}\right)^2\frac{b-u}{u-a}
\]
shows $a<v<b$. It also shows that $u$ and $v$ lie on opposite sides
of $m$, unless both equal $m$. Thus this pointwise construction
defines a strict symmetric mean $V=S_{L_c}(N)$.
If $D=(u-a)(m-b)^2-(m-a)^2(u-b)>0$ is its denominator, polynomial
expansion gives the identity
\begin{equation}\label{certificate}
 D\,Q_c(u,v;m)=-(a-b)(u-m)^2 Q_c(a,b;m)=0.
\end{equation}
For $u\ne v$, $m$ is between $u$ and $v$ and hence is the unique
interior root, giving $L_c(u,v)=m$. For $u=v=m$ the same assertion
is immediate. The diagonal $a=b$ is also immediate.

Finally, a value $w$ satisfying $L_c(u,w)=m$ must satisfy the equation
$Q_c(u,w;m)=0$, which is linear in $w$. Its coefficient is
$2cu-(1+2c)m$. This cannot vanish for a solvable equation: for $c=0$
it equals $-m$; for $c\ne0$, if it vanished the constant term would
be $-m^2/(2c)\ne0$. Therefore $v$ is the unique functional-symmetry
value. This proves the converse and finishes Theorem \ref{main}.

\section*{All asymptotic coefficients}
Writing $L_c(x-t,x+t)=x+h$, \eqref{quad} reduces to
\[
 xh+(1+c)h^2=ct^2.
\]
With $H=h/x$ and $z=t^2/x^2$, this becomes
$H+(1+c)H^2=cz$. The branch $H(0)=0$ is analytic in $z$.
Its coefficients $h=\sum_{n\ge1}b_n t^{2n}x^{1-2n}$ satisfy
\[
 b_1=c,\qquad
 b_n=-(1+c)\sum_{j=1}^{n-1}b_jb_{n-j}\quad(n\ge2).
\]
The Catalan recurrence therefore proves
\[
 b_n=(-1)^{n-1}C_{n-1}c^n(1+c)^{n-1},\qquad
 C_k=\frac1{k+1}\binom{2k}{k}.
\]
This includes $c=-1$ (the harmonic mean) and $c=0$ (the arithmetic
mean), where the higher terms vanish. Since $M=L_c$ globally, its
two-variable expansion also agrees after replacing $x$ by
$x+(s+t)/2$ and the half-difference by $(t-s)/2$.

\section*{Four follow-up tasks}
These are proposed questions, not additional solved credits.
\begin{description}
\item[Task 2.] Classify all continuous strict symmetric homogeneous
means whose two symmetries coincide for every strict symmetric mean,
without assuming any diagonal asymptotic expansion.
\item[Task 3.] Classify all $C^2$ strict symmetric means on
$(0,\infty)^2$ whose two symmetries coincide for every strict symmetric
mean, without assuming homogeneity.
\item[Task 4.] Classify the analytic strict symmetric homogeneous means
$M$ satisfying only $M(A,S_M(A))=M$, where $A$ is the arithmetic mean.
Determine whether this single test forces $M=L_c$.
\item[Task 5.] For fixed $R>1$, $K>0$ and $0<\eta<1/2$, determine the
optimal modulus $\omega(\varepsilon)\to0$, if one exists, bounding
$\inf_c\sup_{[1,R]^2}|M-L_c|$ for $C^4$ symmetric homogeneous means
with $\|M\|_{C^4([1,R]^2)}\le K$, with
$a+\eta(b-a)\le M(a,b)\le b-\eta(b-a)$ for $1\le a<b\le R$,
and with $|M(u,S_m(u))-m|\le\varepsilon$ for all such $a,b$ and
$a<u<b$. Here $m=M(a,b)$ and $S_m(u)$ is \eqref{S}.
\end{description}

\section*{Verification and review}
The proof is analytic and algebraic. The accompanying exact SymPy
verifier checks the rational identities, quadratic endpoint signs,
discriminant, local second-order coefficient and Catalan recurrence.
These checks supplement the proof; they do not replace the limiting
argument or the universal quantifiers. The current 16-page manuscript
and 15-page final article were read completely. Source discrepancies
and the limits of the prior-claim search are in the literature audit.
This is author-agent review, without independent peer review.

\begin{thebibliography}{9}
\raggedright
\bibitem{M} L. Mihokovi\'c, \emph{Coinciding Mean of the Two Symmetries
on the Set of Mean Functions}, Axioms \textbf{12}(3) (2023), 238,
\href{https://doi.org/10.3390/axioms12030238}{doi:10.3390/axioms12030238}.
Current manuscript: \href{https://arxiv.org/abs/2112.07489}{arXiv:2112.07489v3}.
\end{thebibliography}
\end{document}
